\begin{tcolorbox}[
colback=gray!5,
colframe=black,
title=\textbf{Principle of Observational Holonomy},
fonttitle=\bfseries,
breakable]
The Liouvillian connection defines the transport of dynamical
pathways on the manifold of generators. Spectroscopic
measurements, however, do not observe this transport directly.
Instead, they measure only its projection onto the observable
subspace selected by the preparation and detection operators.
Consequently, the experimentally reconstructed holonomy is not
the full Liouville-space holonomy, but the holonomy induced on
the observable subspace. Observation is therefore a geometric
projection, and different spectroscopic protocols reconstruct
different projections of the same underlying response geometry.
The basis-misalignment angle $\theta$ is the local coordinate
parameterizing this observable projection. Reconstructing
$\theta$ from nonlinear spectra is therefore equivalent to
reconstructing the local observational geometry.
\end{tcolorbox}

